%! TEX root = main.tex

\begin{exercise}\label{ex:nt-1}
Compute $\varphi(72)$ from the product formula, and verify it by listing
the residues $1 \le k \le 72$ coprime to $72$.
\end{exercise}

\begin{exercise}\label{ex:nt-2}
Find the inverse of $17$ modulo $64$ using the extended Euclidean
algorithm. Show your back-substitution steps and check the answer.
\end{exercise}

\begin{exercise}\label{ex:nt-3}
Compute $3^{100} \pmod{7}$ using \cref{cor:fermat}, showing the
exponent reduction explicitly.
\end{exercise}

\begin{exercise}\label{ex:nt-4}
Find the last two digits of $3^{2026}$, i.e.\ compute $3^{2026} \pmod{100}$.
\end{exercise}

\begin{exercise}\label{ex:nt-5}
Solve the system $x \equiv 1 \pmod{4}$, $x \equiv 2 \pmod{5}$,
$x \equiv 3 \pmod{7}$, giving the smallest nonnegative solution and
describing all solutions.
\end{exercise}

\begin{exercise}\label{ex:nt-6}
Prove that the CRT system $x \equiv 0 \pmod{6}$, $x \equiv 1 \pmod{10}$
has no solution. Explain which hypothesis of \cref{thm:crt} fails.
\end{exercise}

\begin{exercise}\label{ex:nt-7}
Let $p$ and $q$ be distinct primes and $N = pq$. Prove that
$a^{1 + k\varphi(N)} \equiv a \pmod{N}$ for every integer $a$ and every
$k \ge 0$, including $a$ not coprime to $N$.
\end{exercise}
